\documentclass[bimj,fleqn]{w-art}\usepackage[]{graphicx}\usepackage[]{color}
%% maxwidth is the original width if it is less than linewidth
%% otherwise use linewidth (to make sure the graphics do not exceed the margin)
\makeatletter
\def\maxwidth{ %
  \ifdim\Gin@nat@width>\linewidth
    \linewidth
  \else
    \Gin@nat@width
  \fi
}
\makeatother

\definecolor{fgcolor}{rgb}{0.345, 0.345, 0.345}
\newcommand{\hlnum}[1]{\textcolor[rgb]{0.686,0.059,0.569}{#1}}%
\newcommand{\hlstr}[1]{\textcolor[rgb]{0.192,0.494,0.8}{#1}}%
\newcommand{\hlcom}[1]{\textcolor[rgb]{0.678,0.584,0.686}{\textit{#1}}}%
\newcommand{\hlopt}[1]{\textcolor[rgb]{0,0,0}{#1}}%
\newcommand{\hlstd}[1]{\textcolor[rgb]{0.345,0.345,0.345}{#1}}%
\newcommand{\hlkwa}[1]{\textcolor[rgb]{0.161,0.373,0.58}{\textbf{#1}}}%
\newcommand{\hlkwb}[1]{\textcolor[rgb]{0.69,0.353,0.396}{#1}}%
\newcommand{\hlkwc}[1]{\textcolor[rgb]{0.333,0.667,0.333}{#1}}%
\newcommand{\hlkwd}[1]{\textcolor[rgb]{0.737,0.353,0.396}{\textbf{#1}}}%
\let\hlipl\hlkwb

\usepackage{framed}
\makeatletter
\newenvironment{kframe}{%
 \def\at@end@of@kframe{}%
 \ifinner\ifhmode%
  \def\at@end@of@kframe{\end{minipage}}%
  \begin{minipage}{\columnwidth}%
 \fi\fi%
 \def\FrameCommand##1{\hskip\@totalleftmargin \hskip-\fboxsep
 \colorbox{shadecolor}{##1}\hskip-\fboxsep
     % There is no \\@totalrightmargin, so:
     \hskip-\linewidth \hskip-\@totalleftmargin \hskip\columnwidth}%
 \MakeFramed {\advance\hsize-\width
   \@totalleftmargin\z@ \linewidth\hsize
   \@setminipage}}%
 {\par\unskip\endMakeFramed%
 \at@end@of@kframe}
\makeatother

\definecolor{shadecolor}{rgb}{.97, .97, .97}
\definecolor{messagecolor}{rgb}{0, 0, 0}
\definecolor{warningcolor}{rgb}{1, 0, 1}
\definecolor{errorcolor}{rgb}{1, 0, 0}
\newenvironment{knitrout}{}{} % an empty environment to be redefined in TeX

\usepackage{alltt}
\usepackage{times}
\usepackage{w-thm}
\usepackage[authoryear]{natbib}
\setlength{\bibsep}{2pt}
\setlength{\bibhang}{2em}
\newcommand{\J}{J\"{o}reskog}
  \newcommand{\So}{S\"{o}rbom}
  \newcommand{\bcx}{{\bf X}}
  \newcommand{\bcy}{{\bf Y}}
  \newcommand{\bcz}{{\bf Z}}
  \newcommand{\bcu}{{\bf U}}
  \newcommand{\bcv}{{\bf V}}
  \newcommand{\bcw}{{\bf W}}
  \newcommand{\bci}{{\bf I}}
  \newcommand{\bch}{{\bf H}}
  \newcommand{\bcb}{{\bf B}}
  \newcommand{\bcr}{{\bf R}}
  \newcommand{\bcm}{{\bf M}}
  \newcommand{\bcf}{{\bf F}}
  \newcommand{\bcg}{{\bf G}}
  \newcommand{\bcs}{{\bf S}}
  \newcommand{\bca}{{\bf A}}
  \newcommand{\bcd}{{\bf D}}
  \newcommand{\bcc}{{\bf C}}
  \newcommand{\bce}{{\bf E}}
  \newcommand{\ba}{{\bf a}}
  \newcommand{\bb}{{\bf b}}
  \newcommand{\bc}{{\bf c}}
  \newcommand{\bd}{{\bf d}}
  \newcommand{\bx}{{\bf x}}
  \newcommand{\by}{{\bf y}}
  \newcommand{\bz}{{\bf z}}
  \newcommand{\bu}{{\bf u}}
  \newcommand{\bv}{{\bf v}}
  \newcommand{\bh}{{\bf h}}
  \newcommand{\bl}{{\bf l}}
  \newcommand{\be}{{\bf e}}
  \newcommand{\br}{{\bf r}}
  \newcommand{\bw}{{\bf w}}
  \newcommand{\de}{\stackrel{D}{=}}
  \newcommand{\bt}{\bigtriangleup}
  \newcommand{\bfequiv}{\mbox{\boldmath $\equiv$}}
  \newcommand{\bmu}{\mbox{\boldmath $\mu$}}
  \newcommand{\bnu}{\mbox{\boldmath $\nu$}}
  \newcommand{\bxi}{\mbox{\boldmath $\xi$}}
  \newcommand{\btau}{\mbox{\boldmath $\tau$}}
  \newcommand{\bgamma}{\mbox{\boldmath $\Gamma$}}
  \newcommand{\bphi}{\mbox{\boldmath $\Phi$}}
  \newcommand{\bfphi}{\mbox{\boldmath $\varphi$}}
  \newcommand{\bfeta}{\mbox{\boldmath $\eta$}}
  \newcommand{\bpi}{\mbox{\boldmath $\Pi$}}
  \newcommand{\bequiv}{\mbox{\boldmath $\equiv$}}
  \newcommand{\bvarepsilon}{\mbox{\boldmath $\varepsilon$}}
  \newcommand{\btriangle}{\mbox{\boldmath $\triangle$}}
  \newcommand{\bdelta}{\mbox{\boldmath $\Delta$}}
  \newcommand{\beps}{\mbox{\boldmath $\epsilon$}}
  \newcommand{\btheta}{\mbox{\boldmath $\theta$}}
  \newcommand{\balpha}{\mbox{\boldmath $\alpha$}}
  \newcommand{\bsphi}{\mbox{\boldmath $\varphi$}}
  \newcommand{\bsig}{\mbox{\boldmath $\sigma$}}
  \newcommand{\bfpsi}{\mbox{\boldmath $\psi$}}
  \newcommand{\bfdelta}{\mbox{\boldmath $\delta$}}
  \newcommand{\bsigma}{{\bf \Sigma}}
  \newcommand{\bzero}{{\bf 0}}
  \newcommand{\bpsi}{\mbox{\boldmath $\Psi$}}
  \newcommand{\bep}{\mbox{\boldmath $\epsilon$}}
  \newcommand{\bomega}{\mbox{\boldmath $\Omega$}}
  \newcommand{\bfomega}{\mbox{\boldmath $\omega$}}
  \newcommand{\blambda}{\mbox{\boldmath $\Lambda$}}
  \newcommand{\bflambda}{\mbox{\boldmath $\lambda$}}
  \newcommand{\bfsigma}{\mbox{\boldmath $\sigma$}}
  \newcommand{\bfpi}{{\mbox{\boldmath $\pi$}}}
  \newcommand{\bupsilon}{\mbox{\boldmath $\upsilon$}}
  \newcommand{\obs}{{\rm obs}}
  \newcommand{\mis}{{\rm mis}}
  \theoremstyle{plain}
  \newtheorem{criterion}{Criterion}
  \theoremstyle{definition}
  \newtheorem{condition}[theorem]{Condition}
  \usepackage[]{graphicx}
  \chardef\bslash=`\\ % p. 424, TeXbook
  \newcommand{\ntt}{\normalfont\ttfamily}
  \newcommand{\cn}[1]{{\protect\ntt\bslash#1}}
  \newcommand{\pkg}[1]{{\protect\ntt#1}}
  \let\fn\pkg
  \let\env\pkg
  \let\opt\pkg
  \hfuzz1pc % Don't bother to report overfull boxes if overage is < 1pc
  \newcommand{\envert}[1]{\left\lvert#1\right\rvert}
  \let\abs=\envert
\IfFileExists{upquote.sty}{\usepackage{upquote}}{}
\begin{document}
  %\DOIsuffix{bimj.DOIsuffix}
  \DOIsuffix{bimj.200100000}
  \Volume{xx}
  \Issue{xx}
  \Year{2017}
  \pagespan{1}{}
  \keywords{Censoring by death; Global rank test; Non-inferiority; Key word four; Key word five;\\
  \noindent \hspace*{-4pc} {\small\it (Up to five keywords are allowed and should be given in alphabetical order. Please capitalize the key}\\
  \hspace*{-4pc} {\small\it words)}\\[1pc]
  \noindent\hspace*{-4.2pc} Supporting Information for this article is available from the author or on the WWW under\break \hspace*{-4pc} \underline{http://dx.doi.org/10.1022/bimj.XXXXXXX} (please delete if not
  applicable)
  }  %%% semicolon and fullpoint added here for keyword style

  \title[Running title]{Power of the Wilcoxon-Mann-Whitney test for non-inferiority in the presence
of death-censored observations}
  %% Information for the first author.
  \author[Schmidtmann {\it{et al.}}]{Irene Schmidtmann\footnote{Corresponding author: {\sf{e-mail: Irene.Schmidtmann@uni-mainz.de}}, Phone: +49-6131-173951, Fax: +49-6131-172968}\inst{,1}}
  \address[\inst{1}]{Institute for Medical Biostatistics, Epidemiology and Informatics (IMBEI),
  University Medical Center Johannes Gutenberg University Mainz, D 55101 Mainz}
  %%%%    Information for the second author
  \author[Konstantinides]{Stavros Konstantinides\inst{2}}
  \address[\inst{2}]{Center for Thrombosis and Hemostasis (CTH), University Medical Center Johannes Gutenberg University Mainz, D 55101 Mainz}
  %%%%    Information for the third author
  \author[Binder]{Harald Binder\inst{3}}
  \address[\inst{3}]{Institute for Medical Biometry and Statistics, University Freiburg, Stefan-Meier-Str. 26, D 79104 Freiburg}
  %%%%    \dedicatory{This is a dedicatory.}
  \Receiveddate{zzz} \Reviseddate{zzz} \Accepteddate{zzz}

  \begin{abstract}
  Start. This is the abstract. You may use this LaTeX template to write your text.
  The abstract should not contain multiple paragraphs, formulae NOR references and should ideally
  not be longer than 250 words. This is the abstract. You may use this LaTeX template to write your text.
  The abstract should not contain multiple paragraphs, formulae NOR references and should ideally
  not be longer than 250 words.  This is the abstract. You may use this LaTeX template to write your text.
  The abstract should not contain multiple paragraphs, formulae NOR references and should ideally
  not be longer than 250 words.  This is the abstract. You may use this LaTeX template to write your text.
  The abstract should not contain multiple paragraphs, formulae NOR references and should ideally
  not be longer than 250 words.  This is the abstract. You may use this LaTeX template to write your text.
  The abstract should not contain multiple paragraphs, formulae NOR references and should ideally
  not be longer than 250 words. Start. This is the abstract. You may use this LaTeX template to write your text.
  The abstract should not contain multiple paragraphs, formulae NOR references and should ideally
  not be longer than 250 words.
  Stop.
  \end{abstract}



  %% maketitle must follow the abstract.
  \maketitle                   % Produces the title.

  %% If there is not enough space inside the running head
  %% for all authors including the title you may provide
  %% the leftmark in one of the following three forms:

  %% \renewcommand{\leftmark}
  %% {First Author: A Short Title}

  %% \renewcommand{\leftmark}
  %% {First Author and Second Author: A Short Title}

  %% \renewcommand{\leftmark}
  %% {First Author et al.: A Short Title}

  %% \tableofcontents  % Produces the table of contents.


  \section{Introduction}

  In clinical trials for cardiovascular diseases, often the aim of an
  intervention is to improve function, measured by some quantitative endpoint
  such as the distance walked in the six minute walk test, a biomarker, or ... .
  Typically, unfavourable outcomes of the quantitative endpoint are related to a
  higher risk of cardiovascular death. So, if the patient's condition is severe,
  there is a non-negligible probability of a fatal outcome. Thus, censoring by
  death may occur if a patient dies before the quantitative outcome can be
  determined. Censoring by death leads to missing values which are most unlikely
  to be missing at random. Therefore, ignoring and excluding this kind of
  missing values from analysis does not only decrease power but may also
  introduce bias.

  Worst-rank scores have been considered e. g. by Wittes et al (1989) and Lachin
  (1999) to tackle this problem. Without loss of generality we assume that
  higher values of the quantitative outcome are more favourable than lower
  values. We further assume that any quantitative outcome is better than death
  and that earlier death is worse than later death. This leads to an obvious
  ordering of outcomes which motivates the definition of worst-rank scores.

  Tied worst-rank scores are obtained by allocating a rank score corresponding
  to a single value below the minimum observed value of the quantitative
  endpoint. Untied worst-rank scores can be obtained, if the time of death is
  taken into account. The lowest rank score is then allocated to the patient
  that has died first, subsequent deceased individuals receive ranks according
  to their time of death. After the deceased have bean ranked the subsequent
  ranks are allocated to the surviving patients according to their quantitative
  outcome.

  Lachin (1999) has shown that such an approach is unbiased against a restricted
  one-sided alternative which states that the new treatment is better with
  respect to both, mortality and the quantitative endpoint or at least better
  with respect to one criterion and equal with respect to the other.

  Another view of this approach is to consider it as a global ranking of
  multiple endpoints, e. g. some kind of event - such as death - and some
  quantitative measurement - such as a biomarker. This approached was originally
  suggested by O'Brien (1984 - add reference
  // O’Brien PC. Procedures for comparing samples with multiple endpoints.
     Biometrics. 1984;40:1079 –1087. //) for multiple endpoints and
  applied in cardiology by Felker and Maisel [add reference] and

  Felker GM, Anstrom KJ, Rogers JG. A global ranking approach to end
points in trials of mechanical circulatory support devices. J Card Fail.
2008;14:368 –372.

  Matsouaka and Bentensky (2015 - add reference) have derived power and sample
  size formulae for this case, including as special case a normally distributed
  quantitative endpoint and exponentially distributed times to event.

  We here consider the situation where a new treatment is compared to a standard
  treatment and the intention is to show that the new treatment is non-inferior
  to the standard treatment. As in the


  \section{Thrombolyis application}

  \section{Methods}

  \section{Application}

  \section{Discussion}


  \subsection{Second level heading}
  This is normal body text. Please note that the body text must be divided into numbered sections with
  suitable short verbal titles. Sub-headings are allowed but not sub-subheadings. Please use sentence case in the
  headings.

  This is the body text. Please note that cross-references in the body text should be shown as follows:
  (Miller, 1900), (Miller and Baker, 1900) or if three or more authors (Miller {\it{et al}}., 1900)
  \vspace*{12pt}

  \noindent Bullet lists are not allowed. Always use (i), (ii), etc.
  \vspace*{12pt}

  \noindent Sentences should never start with a symbol.
  \vspace*{12pt}

  \noindent Names of software packages and website addresses should be written in {\tt{Courier new, i.e. Stata, the R package
  MASS, http://www.biometrical-journal.com.}}


  \begin{table}[htb]
  \begin{center}
  \caption{The caption of a table.}
  \begin{tabular}{lll}
  \hline
  Description 1 & Description 2 & Description 3\\
  \hline
  Row 1, Col 1 & Row 1, Col 2 & Row 1, Col 3\\
  Row 2, Col 1 & Row 2, Col 2 & Row 2, Col 3\\
  \hline
  \end{tabular}
  \end{center}
  \end{table}
  \begin{equation}
  \left({\theta^{0}_{i}}\atop{\theta^{1}_{i}}\right) \sim N(\theta,\Sigma),\quad {\mathrm{with}}\
  {{\theta}} = \left({\theta_{0}}\atop{\theta_{1}}\right)\ {\mathrm{and}}\ \Sigma =
  \left(\begin{array}{cc}
  \sigma^{2}_{0} & \rho\sigma_{0}\sigma_{1}\\
  \rho\sigma_{0}\sigma_{2} & \sigma^{2}_{1}
  \end{array}\right).
  \end{equation}

  \noindent This is the body text. Only number equations which are referred to in the text body. If equations
  are numbered, these should be numbered continuously throughout the text. Not section wise! Please
  carefully follow the rules for mathematical expressions in the ``Instructions to Authors''.

  \begin{figure}[htb]
  \begin{center}
%  \includegraphics[bb= 0 0 115 87]{empty.eps}
  \caption{The figure caption ($b_{2}b_{2}a^{n}a^{n}$).}
  \end{center}
  \end{figure}
  \begin{acknowledgement}
  An acknowledgement may be placed at the end of the article.
  \end{acknowledgement}
  \vspace*{1pc}

  \noindent {\bf{Conflict of Interest}}

  \noindent {\it{The authors have declared no conflict of interest. (or please state any conflicts of interest)}}

  \section*{Appendix {\it(please insert here, if applicable)}}

  \subsection*{A.1.\enspace Second level heading}

  Please insert appendices before the references.

  \begin{thebibliography}{10}
  \bibitem[Lachin(1999)Lachin, J. L]{bib1}Lachin, J. L. (1999) Worst-rank score analysis with informatively missing observations in clinical trials.  \textit{Controlled Clinical Trials} \textbf{20}, 408–422.

  \bibitem[Wittes et al(1989)Wittes, J., Lakatos, E., and Probstfield, J.]{bib2}Wittes, J., Lakatos, E., and Probstfield, J. (1989) Surrogate endpoints in clinical trials: cardiovascular diseases. \textit{Statistics in Medicine} \textbf{8}, 415–425.


%\bibitem[Bauer and Bauer(1994)Bauer, P. and Bauer, M.M.]{bib1}Bauer, P. and  Bauer, M. M. (1994). Testing equivalence simultaneously for location and  dispersion of two normally distributed populations.  \textit{Biometrical  Journal} \textbf{36}, 643--660.
%  \bibitem[Farrington, C. P. and Andrews, N. (2003)]{bib2}Farrington, C.P. and Andrews, N. (2003). Outbreak detection:   Application to infectious disease surveillance. In: Monitoring the Health of Populations (eds. R. Brookmeyer and D. F. Stroup), Oxford University Press, Oxford,\break 203--231.
  \bibitem[Rencher(1998)Rencher, A.C.]{bib3}Rencher, A. C. (1998).  \textit{Multivariate Statistical Inference and Applications}. Wiley, New  York.
  \end{thebibliography}
  \newpage
  \phantom{aaaa}
  \end{document}
